\documentclass[11pt]{article}
\usepackage[T1]{fontenc}
\usepackage{lmodern}
\usepackage[margin=1in]{geometry}
\usepackage{amsmath,amssymb,amsthm,booktabs,microtype}
\usepackage[hidelinks]{hyperref}
\newtheorem{theorem}{Theorem}
\newtheorem{lemma}[theorem]{Lemma}
\newcommand{\lr}{\mathrel{\leq_{\mathrm{lr}}}}
\newcommand{\supp}{\operatorname{supp}}
\newcommand{\doi}[1]{\href{https://doi.org/#1}{\nolinkurl{#1}}}
\title{A two-step bound for the modes of\ uniform hypercomb summands}
\author{Xiaojian Zeng\thanks{PBC School of Finance, Tsinghua University, 100190.\newline
Email: \href{mailto:xiaojianzeng@foxmail.com}{xiaojianzeng@foxmail.com}.}}
\date{}

\begin{document}
\maketitle
\begin{abstract}
Galvin and Sharpe express the strong independence polynomial of a uniform
linear hypercomb as a sum of two log-concave polynomials. Their mode estimates
leave a possible interval of nonmonotonicity of length $O(\sqrt n)$ for each
fixed uniformity. For every spine length $n\geq2$ and uniformity $\ell\geq3$,
we prove that the smallest modes of these two summands differ by at most two.
More precisely, the summands $W_n,Y_n$ satisfy
$W_n\leq_{\mathrm{lr}}x^2Y_n$ and $Y_n\leq_{\mathrm{lr}}xW_n$ in the
likelihood-ratio order, with coefficients extended by zero.
Their sum is therefore nondecreasing up to the smaller mode and
nonincreasing from the larger. The separation two is attained at
$(n,\ell)=(3,5)$. The interval is located by the summand modes; full
unimodality of the hypercomb polynomial remains unresolved here.
\end{abstract}

\section{The family and the result}

Fix integers $n\geq1$ and $\ell\geq3$. A linear $\ell$-uniform hypercomb
has distinct attachment vertices $w_0,\ldots,w_n$. For $1\leq i\leq n$,
its $i$th spine edge consists of $w_{i-1},w_i$ and $\ell-2$ private vertices.
At each $w_i$ attach a tooth consisting of $w_i$ and $\ell-1$ private
vertices. All private vertex sets are disjoint. In particular, the two
attachments are distinct when $n=1$. A \emph{strong independent set}
meets each hyperedge in at most one vertex. Write
\[
 C_n(x)=C_{n,\ell}(x)=\sum_{k=0}^{2n+1}c_{n,\ell,k}x^k
\]
for its counting polynomial. Each edge has a private vertex, so the maximum
size of such a set is $2n+1$.

Galvin and Sharpe~\cite[Theorem 1.7, Corollary 1.8 and Section 2.5]{GS}
give a recurrence for $C_n$ and a decomposition into two log-concave
summands. For each fixed $\ell$, their bounds give monotonicity outside
an interval $[K_\ell n-L_\ell\sqrt n,K_\ell n+L_\ell\sqrt n]$.
Their Question 1.10(4) asks whether this interval can be shortened or
eliminated. Note 1.9 of that paper withdraws an earlier proof of full
log-concavity. The theorem below shortens the possible interval, with
endpoints specified by summand modes rather than an explicit affine center.

Put $s=\ell-2$ and $t=\ell-1=s+1$, and define
\begin{equation}\label{eq:factors}
 \begin{aligned}
 a&=(1+sx)(1+tx), & b&=x(1+tx),\\
 z_0&=1+(t+1)x, & z_1&=(1+tx)(a+2x).
 \end{aligned}
\end{equation}
The recurrence in~\cite{GS} reads
\begin{equation}\label{eq:comb}
 C_0=z_0,\qquad C_1=z_1,\qquad C_n=aC_{n-1}+bC_{n-2}\quad(n\geq2).
\end{equation}
Here $C_0$ is an auxiliary initial polynomial. Let
\begin{equation}\label{eq:F}
 F_0=0,\qquad F_1=1,\qquad F_j=aF_{j-1}+bF_{j-2}\quad(j\geq2).
\end{equation}
Then, by induction in~\eqref{eq:comb}, the summands from
\cite[Section 2.5, equations (28)--(30)]{GS} are
\begin{equation}\label{eq:split}
 C_n=W_n+Y_n,\qquad W_n=z_1F_n,\qquad Y_n=z_0bF_{n-1}\quad(n\geq2).
\end{equation}

For nonzero polynomials $P=\sum_i p_ix^i$ and $Q=\sum_i q_ix^i$ with
nonnegative coefficients, extended by zero to all integer indices, set
\begin{equation}\label{eq:order}
 P\lr Q\quad\Longleftrightarrow\quad
 p_iq_j-p_jq_i\geq0\quad\text{for every }i<j.
\end{equation}
This is the usual likelihood-ratio order without normalization; see
\cite[Section 2]{DM}. Write $m(P)=\min\{i:p_i=\max_j p_j\}$ for the
\emph{smallest} mode.

\begin{theorem}\label{thm:main}
For every $n\geq2$ and $\ell\geq3$, the summands in~\eqref{eq:split} satisfy
\begin{equation}\label{eq:mainorder}
 W_n\lr x^2Y_n,\qquad Y_n\lr xW_n.
\end{equation}
Consequently,
\begin{equation}\label{eq:modes}
 m(Y_n)-1\leq m(W_n)\leq m(Y_n)+2.
\end{equation}
If $r=\min\{m(W_n),m(Y_n)\}$ and $q=\max\{m(W_n),m(Y_n)\}$, then
$q-r\leq2$ and
\begin{equation}\label{eq:tails}
 c_{n,\ell,0}\leq\cdots\leq c_{n,\ell,r},\qquad
 c_{n,\ell,q}\geq\cdots\geq c_{n,\ell,2n+1}.
\end{equation}
The separation bound two is attained by unique summand modes $6$ and $4$
at $n=3$, $\ell=5$.
\end{theorem}

The contribution is the simultaneous order invariant
$F_{j-1}\lr F_j\lr bF_{j-1}$ and its application to the initial factors.
The recurrence, factorization and summand log-concavity are prior ingredients.
Recurrence factorizations for graph independence polynomials occur in
Wang and Zhu~\cite[Theorem 3.1]{WZ}; related order-based arguments appear
in Bautista-Ramos, Guill\'en-Galv\'an and G\'omez-Salgado~\cite{BGG}.
Their corrigendum~\cite{BGGcorr} addresses support conditions in ratio
comparisons. We retain all pairs of indices in~\eqref{eq:order} and prove
the required order properties, including zero coefficients, below.

The September 2026 preprint~\cite[Theorem 1 and Corollary 2]{Eventual}
asserts eventual strict log-concavity for each fixed pair of positive
weights. Its recurrence specializes to~\eqref{eq:comb} at
$s=\ell-2$, $t=\ell-1$; its Section 9 retains the question at every
length in the uniform family. Theorem~\ref{thm:main} applies at every
length in its stated range and uses no result from that preprint.
Its analytic proof and reported formalization are not evaluated here.
The comparison with these sources, checked on 14 September 2026, is a
bounded assessment of the contribution, not a claim of worldwide priority.

\section{Order calculus and log-concavity}

In this section every nonzero sequence has finite support and is positive
at every integer between its support endpoints. A polynomial is
\emph{log-concave} if $p_i^2\geq p_{i-1}p_{i+1}$ for all $i$.
With the support convention above, log-concavity implies unimodality.

\begin{lemma}\label{lem:order}
The order~\eqref{eq:order} is transitive. Its inequalities are preserved
under addition with a common lower or upper bound, and under multiplication
by a log-concave polynomial with interval support. Moreover,
$P\lr Q$ implies $m(P)\leq m(Q)$.
\end{lemma}
\begin{proof}
If $P\lr Q$, each support endpoint of $P$ is at most the corresponding
endpoint of $Q$. For example, if the left endpoint of $Q$ preceded that
of $P$, the minor on these two endpoints would be negative; the right
endpoints are treated similarly. Suppose $P\lr Q\lr R$ and $i<j$.
If $p_jr_i=0$, the required inequality is immediate. Otherwise the
endpoint inequalities place $i,j$ in the positive support of $Q$.
Multiplying $p_iq_j\geq p_jq_i$ and $q_ir_j\geq q_jr_i$ and dividing
by $q_iq_j>0$ proves transitivity. Linearity of each minor in either
polynomial proves the assertions about addition whenever the resulting
sequences have interval support, as they do below.

Let $R=\sum_k r_kx^k$ be log-concave. Decreasing successive coefficient
ratios give
\begin{equation}\label{eq:toeplitz}
 r_{i-u}r_{j-v}-r_{j-u}r_{i-v}\geq0\qquad(i<j,\ u<v).
\end{equation}
Indeed, if the second product is positive, its two indices are the
outermost of the four indices. The indices of the first product lie
between them and have the same sum. All four coefficients are then
positive, and the ratio argument applies. If the second product vanishes,
the inequality is immediate. Expanding the convolution minor gives
\begin{equation}\label{eq:convolution}
 \begin{aligned}
 &(PR)_i(QR)_j-(PR)_j(QR)_i\\
 &\quad=\sum_{u<v}(p_uq_v-p_vq_u)
       (r_{i-u}r_{j-v}-r_{j-u}r_{i-v})\geq0.
 \end{aligned}
\end{equation}
Here $(PR)_i$ denotes the coefficient of $x^i$ in $PR$. This proves
preservation under multiplication.

Finally, if $i=m(Q)<j=m(P)$, then $p_i<p_j$, $q_i\geq q_j$ and
$q_i>0$. Thus $p_iq_j<p_jq_i$, contradicting~\eqref{eq:order}.
This also covers tied modes.
\end{proof}

We will use the classical closure of log-concavity under convolution,
stated, for example, in~\cite[Theorem 1.4]{JG}. It follows directly from
the lemma: log-concavity and interval support imply $P\lr xP$ by the
ratio argument. Multiplication by a log-concave $R$ gives
$PR\lr xPR$, whose adjacent minors are precisely the log-concavity
inequalities for $PR$. Products preserve interval support.

\begin{lemma}\label{lem:lc}
The polynomials $a,b,z_0,z_1$, and every $F_j$ with $j\geq1$, are
log-concave. The supports of $F_j,W_n,Y_n$ are respectively
$[0,2j-2]$, $[0,2n+1]$ and $[1,2n-1]$.
\end{lemma}
\begin{proof}
The factors $a,b,z_0$ are products of linear polynomials and monomials.
Also $d=a+2x=1+(s+t+2)x+stx^2$ is log-concave since
$(s+t+2)^2>st$. Hence $z_1=(1+tx)d$ is log-concave.

For $0\leq y\leq1$, factor
\begin{equation}\label{eq:cubic}
 \begin{aligned}
 a^2+4yb&=(1+tx)R_y,\\
 R_y&=1+(2s+t+4y)x+(s^2+2st)x^2+s^2tx^3.
 \end{aligned}
\end{equation}
The two interior log-concavity determinants of $R_y$ are positive:
\begin{align}
 (2s+t+4y)^2-(s^2+2st)
   &=3s^2+2st+t^2+16sy+8ty+16y^2>0,\label{eq:det1}\\
 (s^2+2st)^2-(2s+t+4y)s^2t
   &=s^2(s^2+2st+3t^2-4ty)\notag\\
   &\geq s^2(6s^2+4s-1)>0.\label{eq:det2}
\end{align}
The last inequality uses $t=s+1$, $s\geq1$ and $y\leq1$.
Thus $a^2+4yb$ is log-concave. These identities give an algebraic
verification of the factor condition in~\cite[Section 2.5]{GS}.

The Fibonacci recurrence~\eqref{eq:F} has the factorization
\begin{equation}\label{eq:factorization}
 F_j=a^{\varepsilon_j}
 \prod_{h=1}^{\lfloor(j-1)/2\rfloor}
 \left(a^2+4b\cos^2\frac{h\pi}{j}\right),
 \qquad
 \varepsilon_j=\begin{cases}1,&j\text{ even},\\0,&j\text{ odd}.
 \end{cases}
\end{equation}
An empty product equals one. To see the identity, solve the recurrence as
$F_j=(\lambda^j-\mu^j)/(\lambda-\mu)$, where
$\lambda+\mu=a$ and $\lambda\mu=-b$. Pairing conjugate factors of the
numerator yields~\eqref{eq:factorization}, with the unpaired factor
$\lambda+\mu$ when $j$ is even. One may first take $x>0$, where the
characteristic roots are distinct, and then use polynomial identity.
This is the factorization underlying~\cite[equations (29)--(30)]{GS}.

Each cosine square lies in $[0,1]$, so product closure proves
log-concavity of $F_j$. Since $a$ has positive coefficients in degrees
$0,1,2$, induction in~\eqref{eq:F} gives the support $[0,2j-2]$.
Multiplication by $z_1$ and $z_0b$ gives the asserted supports and
log-concavity of the summands.
\end{proof}

\section{The invariant and the two-step bound}

\begin{lemma}\label{lem:invariant}
For every $j\geq2$,
\begin{equation}\label{eq:invariant}
 F_{j-1}\lr F_j\lr bF_{j-1}.
\end{equation}
\end{lemma}
\begin{proof}
The base case is $1=F_1\lr F_2=a\lr b$. The only potentially
nontrivial minors of $a\lr b$, at $(0,1),(0,2),(1,2)$, are
$1,t,t^2$, respectively.

Assume~\eqref{eq:invariant}. Because $1\lr a$ and $F_j$ is
log-concave, Lemma~\ref{lem:order} gives $F_j\lr aF_j$.
Together with the right inequality in~\eqref{eq:invariant}, addition
gives
\[
 F_j\lr aF_j+bF_{j-1}=F_{j+1}.
\]
For the other half, $a\lr b$ and log-concavity of $F_j$ give
$aF_j\lr bF_j$. The left inequality in~\eqref{eq:invariant},
multiplied by the log-concave kernel $b$, gives
$bF_{j-1}\lr bF_j$. Adding the lower polynomials yields
$F_{j+1}\lr bF_j$, completing the simultaneous induction.
\end{proof}

\begin{proof}[Proof of Theorem~\ref{thm:main}]
Expand
\[
 z_1=1+(3\ell-2)x+3t^2x^2+st^2x^3.
\]
First, $z_1\lr x^2z_0$: the only overlap pair needing comparison is
$(2,3)$, with minor $t^2(3(t+1)-s)>0$. Every other minor is
nonnegative by the supports. Apply Lemma~\ref{lem:invariant} and this
initial-factor comparison, with log-concave kernels $z_1$ and
$bF_{n-1}$, respectively:
\begin{equation}\label{eq:chain1}
 W_n=z_1F_n\lr z_1bF_{n-1}\lr x^2z_0bF_{n-1}=x^2Y_n.
\end{equation}
Second, $z_0\lr d=a+2x$, since its minors at
$(0,1),(0,2),(1,2)$ are $s+1,st,(t+1)st$. Multiplication by
$b=x(1+tx)$ gives $z_0b\lr xz_1$. Using the log-concave kernels
$z_0b$ and $F_n$ gives
\begin{equation}\label{eq:chain2}
 Y_n=z_0bF_{n-1}\lr z_0bF_n\lr xz_1F_n=xW_n.
\end{equation}
Transitivity proves~\eqref{eq:mainorder}.

Lemma~\ref{lem:order} and $m(x^hP)=m(P)+h$ now give~\eqref{eq:modes}.
Both summands are log-concave by Lemma~\ref{lem:lc}. Adding their
monotonicity inequalities proves~\eqref{eq:tails}. Finally, the exact
coefficients in Table~\ref{tab:witness} show the unique modes $6$ and
$4$, so the bound on summand separation is attained.
\end{proof}

\begin{table}[htbp]
\centering
\caption{Exact coefficients at $n=3$, $\ell=5$. Omitted high-degree
coefficients of $Y_n$ are zero.}\label{tab:witness}
\small
\begin{tabular}{crrrrrrrr}
\toprule
$k$ & 0 & 1 & 2 & 3 & 4 & 5 & 6 & 7\\
\midrule
$[x^k]W_n$ & 1 & 28 & 320 & 1937 & 6744 & 13632 & 14976 & 6912\\
$[x^k]Y_n$ & 0 & 1 & 16 & 95 & 248 & 240 & 0 & 0\\
$[x^k]C_n$ & 1 & 29 & 336 & 2032 & 6992 & 13872 & 14976 & 6912\\
\bottomrule
\end{tabular}
\end{table}

The sum in Table~\ref{tab:witness} is itself log-concave. Thus the example
establishes sharpness only for the separation of these summands. The case
$n=1$ is log-concave directly from $C_1=z_1$. In general, when $q-r\leq1$,
the two monotone tails already give unimodality. When $q=r+2$, the only
remaining possible failure is the strict valley
\[
 c_{n,\ell,r}>c_{n,\ell,r+1}<c_{n,\ell,r+2}.
\]
The argument does not exclude this valley. Neither full unimodality nor
log-concavity for all parameters follows from the theorem.

\section{Exact checks and reproducibility}

The infinite conclusion follows from the preceding proof. Finite exact
checks give a separate test of the formulas and their implementation.
An independently written verifier generated $F_j$ by the tiling expansion
\[
 F_j=\sum_{h=0}^{\lfloor(j-1)/2\rfloor}
       \binom{j-1-h}{h}a^{j-1-2h}b^h
\]
and tested the recurrence, all four order comparisons in
\eqref{eq:invariant} and~\eqref{eq:mainorder}, the summand
log-concavity and the two monotone tails. It checked all $i<j$ minors
for all pairs in
\[
 \bigl(\{2,\ldots,24\}\times\{3,\ldots,18\}\bigr)
 \ \cup\
 \bigl(\{25,32,47,64\}\times\{3,4,5,31,1000003\}\bigr).
\]
These are $388$ parameter pairs and $958{,}808$ minors, all passing.
A second construction enumerated subsets $A\subseteq\{0,\ldots,n\}$
of occupied attachment vertices with no consecutive positions. If
$h(A)$ counts spine edges with neither endpoint in $A$, the contribution is
\[
 x^{|A|}(1+tx)^{n+1-|A|}(1+sx)^{h(A)}.
\]
Summing these contributions agreed with $C_n$ for $1\leq n\leq10$,
$3\leq\ell\leq8$, giving $60$ direct combinatorial comparisons.

Independent reconstruction is also used in exact verification for related
pendant-edge constructions: Zeng~\cite{ZengWhiskered} reports integer
kernel certificates, checkable by reconstructing multiplication matrices,
for the eight-vertex whiskered-graph weak Lefschetz problem. Here the
reconstructed objects are coefficient arrays and hypercomb counts.
The accompanying verifier uses Python integers and no external package;
it can be run as \texttt{python3 manuscript/verify.py} from the project
directory. It recomputes the finite checks just described without modifying
the research records. The finite ranges impose no cutoff on
Theorem~\ref{thm:main}.

\paragraph{Computational and AI assistance.}
OpenAI Codex was used in the mathematical exploration, proof development,
literature searches, exact-check programs and manuscript preparation.
A separate model context reviewed the mathematical argument and wrote the
second verifier. This was an additional review pass, not external peer review
or statistically independent evidence. The initial symbolic checks used
SymPy 1.14.0; the accompanying verification uses exact integer arithmetic.
No formal proof-assistant verification of this note is claimed.

\bibliographystyle{plain}
\bibliography{references}
\end{document}
